\chapter{Floating-Point Arithmetic and Error Analysis}

A program solves $x^{2}+10^{8}x+1=0$ with the quadratic formula we have used since
school, $x=(-b\pm\sqrt{b^{2}-4ac})/(2a)$. For $a=1$, $b=10^{8}$, $c=1$ it reports
the roots $-10^{8}$ and $-7.450580596923828\times10^{-9}$. The first root is correct
to the last digit. The second is wrong by a quarter: the true value is $-10^{-8}$, and
its relative error is $0.255$. The formula is not mistyped. Neither is
the machine broken. The coefficients are exact, the discriminant is exact, and the
subtraction is even exact according to the theorem we shall prove below.

What collapses is the algebraic form. The numbers $b=10^{8}$ and $\sqrt{b^{2}-4ac}$
differ by only one ulp, that is one smallest step the machine can still tell apart at
that scale: the difference is $1.49\times10^{-8}$, or $1.5\times10^{-16}$ times the
size of the numbers themselves. A difference that small means nothing as long as we
regard numbers of size $10^{8}$. As soon as the two are subtracted, the result has
size $10^{-8}$ and that difference becomes the whole answer: what remains is error,
not a number.

Chapters 4 and 5 build limits, continuity, and derivatives over $\mathbb{R}$ whose
every point is exact. Chapter 20 uses the variance formula and the square root under
the same assumption. A computer does not work in $\mathbb{R}$. It works in a finite
set of numbers that has rules of its own: every operation is rounded, addition is not
associative, and two algebraically equivalent routes can give answers far apart. This
chapter builds a way of reading those rules, of measuring the damage before the program
runs, and of choosing a form of the formula that does not collapse.

\begin{untukanak}
Imagine a digital kitchen scale whose display holds only five digits. We weigh a sack
of rice together with its container, then weigh the empty container, then subtract the
two numbers to get the weight of the rice. If container and rice each weigh about
twenty kilograms, the display can only show numbers to the nearest tens of grams. The
difference of two large numbers that are nearly equal is correct as arithmetic, but its
last digits no longer mean anything. Rice that actually weighs eight hundred grams can
show up as seven hundred grams, or even as an absurd number.

There is a second, subtler problem. The same scale cannot weigh a single grain of rice:
the needle does not move at all. If we pour in a million grains one at a time, not one
of them is recorded, even though a million grains have weight. A calculating machine
works in exactly that way: some numbers are too small to change anything, and some
subtractions throw away almost every digit.

The question: if the machine is like that, how can we know when an answer can still be
trusted? In this chapter we shall learn why that error is not a random event but
something whose bound can be computed before a single line of the program runs.
\end{untukanak}

\section{Numbers without a representative}
\tingkat{basic}

Between $0$ and $1$ there are infinitely many real numbers. Inside a computer, all
numbers are stored in 32 or 64 bits, so only finitely many are available. The
compromise adopted from the first machines down to the IEEE-754 standard is binary
scientific notation with fixed field widths: one sign, one exponent, one mantissa.

\kamus{floating-point system}{sistem floating-point}{Binary scientific notation with a
fixed number of digits, like writing $1.01\times2^{7}$ in limited room.}

\begin{definisi}[Floating-point number system]
Let $\beta=2$ be the base, $p$ the precision, and $e_{\min}\le e\le e_{\max}$ the
exponent bounds. A normal number is
\[
  x=\pm\left(1.d_{1}d_{2}\dots d_{p-1}\right)_{2}\times2^{e},
\]
with $d_{i}\in\{0,1\}$. Subnormal numbers use $e=e_{\min}$ with leading digit
$0$.\index{subnormal}\index{floating-point arithmetic} We call the set of all
representable numbers, together with zero, the two infinities, and NaN,
$\mathbb{F}$.
\end{definisi}

The leading digit is always $1$ for a normal number, so that digit is not stored;
that is why the mantissa field has $p-1$ bits. One bit saved doubles the precision:
23 stored bits give 24 significant
bits.



\begin{teorema}[Number of machine numbers]
For every exponent $e$, the number of normal numbers in the interval
$\left[2^{e},2^{e+1}\right)$ is $2^{p-1}$. The set $\mathbb{F}$ is finite, hence
countable.
\end{teorema}

\begin{proof}
On that interval every number has the form $M\cdot2^{e-p+1}$ with $M$ an integer
and $2^{p-1}\le M<2^{p}$. The number of such $M$ is exactly $2^{p}$, minus
$2^{p-1}$, that is $2^{p-1}$. There are finitely many allowed exponents, so
$\mathbb{F}$ is a finite union of finite sets.
\end{proof}

For binary64, $2^{p-1}=2^{52}=4\,503\,599\,627\,370\,496$. That is a lot, but Chapter 2
already showed that $\mathbb{R}$ is uncountable. So almost every real number has no
representative: representation always means rounding, and rounding always means
error. The question is not whether there is error, but how large it is and where it
goes.

\begin{bacakeras}
Two voices, read alternately with a friend. Tap the table every time the word ``piece''
is said.

``One divided by three, how much?'' ``Zero point three three three, on and on, without
end.'' ``Does the machine write it out whole?'' ``No. It stops where there is room.''
``So what is stored?'' ``A piece. Not the value, the piece.''
``And if the pieces are added?'' ``What grows is the pieces, not the value.''
``And if the pieces are subtracted?'' ``The leading digits are used up, dust is what
remains.''

Repeat the three middle lines, faster and faster, with the same tapping: a piece is not
the value, a piece is not the value, a piece is not the value.

The second voice closes quietly: the machine does not lie, it only cuts. We are the ones
who go wrong if we read the piece as the value.

The line the reader must repeat, once, in a clear voice: a computer stores pieces, not
values, so every computed result must be checked, not trusted.
\end{bacakeras}

The IEEE-754 standard fixes two widths that are used everywhere. This table contains all
the numbers; nothing needs memorizing except $p$ and the bias.

\begin{center}
\begin{tabular}{lcc}
\hline
 & binary32 & binary64 \\
\hline
sign / exponent / mantissa bits & $1/8/23$ & $1/11/52$ \\
precision $p$ & $24$ & $53$ \\
exponent bias & $127$ & $1023$ \\
$e_{\min},\ e_{\max}$ & $-126,\ 127$ & $-1022,\ 1023$ \\
$\mathrm{eps}=2^{-p+1}$ & $1.1920929\times10^{-7}$ & $2.2204460\times10^{-16}$ \\
$u=2^{-p}$ & $5.9604645\times10^{-8}$ & $1.1102230\times10^{-16}$ \\
smallest normal & $1.1754944\times10^{-38}$ & $2.2250739\times10^{-308}$ \\
smallest subnormal & $1.4012985\times10^{-45}$ & $4.9406565\times10^{-324}$ \\
largest & $3.4028235\times10^{38}$ & $1.7976931\times10^{308}$ \\
\hline
\end{tabular}
\end{center}

Exponents are stored with a bias so that no separate sign is needed: a stored value of
$123$ means an actual exponent of $123-127=-4$. This trick makes the bit patterns from
zero up to positive infinity increase monotonically, so that two positive numbers can be
compared as integers.

Let us take apart three bit patterns that Python actually produces. In binary32,
$0.1$ prints as \texttt{0 01111011 10011001100110011001101} with
$e=123-127=-4$; $6$ becomes \texttt{0 10000001 10000000000000000000000}, that is
$1.5\times2^{2}$; and $-12.375$ becomes
\texttt{1 10000010 10001100000000000000000}, that is $-1.546875\times2^{3}$.
Two boundary cases complete the picture:

\begin{center}
\begin{tabular}{ll}
\hline
smallest subnormal & \texttt{0 00000000 00000000000000000000001} \\
largest & \texttt{0 11111110 11111111111111111111111} \\
\hline
\end{tabular}
\end{center}

The first has the value $2^{-149}$: all its exponent bits are zero and only the last
mantissa bit is on. The second has the value $3.4028235\times10^{38}$, with stored
exponent $254$ and a full mantissa.

A more readable form for binary64 is provided by \texttt{float.hex()}. Here is the result for several numbers:

\begin{center}
\begin{tabular}{ll}
\hline
$0.1$ & \texttt{0x1.999999999999ap-4} \\
$6.0$ & \texttt{0x1.8000000000000p+2} \\
$10^{-8}$ & \texttt{0x1.5798ee2308c3ap-27} \\
$5\times10^{-324}$ & \texttt{0x0.0000000000001p-1022} \\
\hline
\end{tabular}
\end{center}

That form shows three things. The digits after the hexadecimal point are always
thirteen, and $13\times4=52$, the width of the binary64 mantissa. The last digit of
$0.1$ is \texttt{a}, that is $1010$, a rounding up from the repeating pattern
$1001$. The last row has a zero before the point: the mark of a subnormal number.

Special values complete this system: $\pm\infty$ for results beyond
$e_{\max}$, signed zero, subnormal numbers below $2^{e_{\min}}$, and NaN
(\emph{not a number}) for indeterminate forms such as $0/0$. One property of NaN
deserves remembering because it often surprises: it is not equal to itself, so that
the comparison \texttt{x == x} is false precisely when $x$ is NaN.\index{NaN}

\section{The only bridge: rounding}
\tingkat{basic}

No operation other than rounding connects $\mathbb{R}$ with $\mathbb{F}$. That is why
the whole of error analysis rests on one question: how far is the nearest machine number
from the original number, and how does that distance scale.

IEEE-754 lays down the rule \emph{round to nearest, ties to even}: choose the nearest
machine number, and if two candidates are equally far, choose the one whose last
mantissa digit is even. The second rule is no theoretical luxury. Without that rule,
rounding at a tie point always leans one way, and the error of a million operations
piles up in the same direction instead of canceling out.

Tie points really do occur, and we can create one deliberately. In
binary64, the distance between $1$ and the number after it is $2^{-52}$. Hence
$1+2^{-53}$ falls exactly halfway, and the even rule chooses $1$ because its mantissa
ends in zero. Python gives

\begin{center}
\texttt{1 + 2.0**-53 = 1.0}, whereas \texttt{1 + 3*2.0**-53 = 1.0000000000000004}.
\end{center}

The first case rounds down, the second up, although both are exactly halfway. Only the
parity of the last digit distinguishes them. This is why
$1+10^{-16}=1$ in binary64: not because $10^{-16}$ is ignored carelessly,
but because there is no room to store it.

\kamus{machine epsilon}{epsilon mesin}{The distance between $1$ and the next machine
number: the size of one step at scale one.}\index{machine epsilon}

Two numbers are often confused, and the confusion changes an error bound by a factor of
two. The machine epsilon $\mathrm{eps}=2^{-p+1}$ is the distance between $1$ and the
number after it, namely $2^{-52}$ for binary64. The unit roundoff
$u=2^{-p}=2^{-53}$ is half that distance, and only $u$ appears in error
bounds. Python calls the first \texttt{numpy.finfo(float).eps}; the second we must
compute ourselves.

The distance between numbers is not uniform. Near $1$ it is $2^{-52}$, near
$10^{10}$ it is $1.9\times10^{-6}$. What is uniform is that distance relative
to the size of the number, and that is the content of the following lemma.

\begin{lemma}[Ulp]
If $2^{e}\le|x|<2^{e+1}$, then the distance between $x$ and the next machine number is
$\operatorname{ulp}(x)=2^{e-p+1}$.
\end{lemma}

\begin{proof}
On that interval all normal numbers have the form $M\cdot2^{e-p+1}$ with $M$ an integer
and $2^{p-1}\le M<2^{p}$; the number of such $M$ is $2^{p-1}$ and the gap between two
consecutive $M$ is one unit $2^{e-p+1}$. The next interval uses $e+1$ so that the gap
doubles, and that is the only point where the density changes.
\end{proof}

\kamus{ulp}{unit in the last place}{The smallest step the machine can tell apart at the
scale of a number.}\index{ulp}

The lemma explains two things at once.
$\operatorname{ulp}(1)=2^{-52}=2.220446049250313\times10^{-16}$, whereas
$\operatorname{ulp}(10^{10})=2^{-52}\cdot2^{33}\approx1.9\times10^{-6}$. So
adding $1$ to $10^{16}$ changes nothing in binary64, because
$\operatorname{ulp}(10^{16})=2$: there really is no machine number between the two.

Figure~\ref{fig:titikmesin} shows the consequence of that lemma over the first four
octaves: each octave holds $2^{23}$ points, but its width doubles, so that the
distance between points doubles as well.

\begin{figure}[htbp]
\centering
\begin{tikzpicture}[x=2.1cm,y=1cm,every node/.style={font=\footnotesize}]
  \draw[ink,thick,->] (-0.55,0) -- (4.30,0);
  \draw[ink,thick,fill=white] (-0.42,0) circle (1.7pt);
  \draw[ink,thick] (-0.48,-0.09) -- (-0.36,0.09);
  \foreach \x in {0,0.125,...,0.875} {\draw[ink] (\x,0.055) -- (\x,-0.055);}
  \foreach \x in {1,1.25,...,1.875} {\draw[ink] (\x,0.055) -- (\x,-0.055);}
  \foreach \x in {2,2.5,...,3.0} {\draw[ink] (\x,0.055) -- (\x,-0.055);}
  \foreach \x in {3,4} {\draw[ink] (\x,0.055) -- (\x,-0.055);}
  \foreach \x/\lab in {0/{$2^{0}$},1/{$2^{1}$},2/{$2^{2}$},3/{$2^{3}$},4/{$2^{4}$}} {
    \draw[ink,thick] (\x,0.14) -- (\x,-0.14);
    \node[below=1pt] at (\x,-0.14) {\lab};
  }
  \draw[{Latex}-{Latex},accent] (0.125,-0.62) -- (0.25,-0.62) node[midway,below] {$d$};
  \draw[{Latex}-{Latex},accent] (1.25,-0.62) -- (1.50,-0.62) node[midway,below] {$2d$};
  \draw[{Latex}-{Latex},accent] (2.50,-0.62) -- (3.00,-0.62) node[midway,below] {$4d$};
  \draw[{Latex}-{Latex},accent] (3.00,-0.62) -- (4.00,-0.62) node[midway,below] {$8d$};
  \node[align=center,text width=3.0cm] at (1.95,0.92)
    {each octave holds $2^{23}$ points, twice the width of the one before};
  \node[align=center,text width=3.0cm] at (0.40,-1.32)
    {leftward toward zero: the gaps shrink};
\end{tikzpicture}
\caption{The first four octaves of binary32 machine numbers. Each octave holds $2^{23}$
points, but its width is twice that of the octave before, so the spacing between points
doubles: $d=2^{-23}$ on $[1,2)$, $2d$ on $[2,4)$, $4d$ on $[4,8)$, and $8d$ on
$[8,16)$. The tail does not hold fewer points, only sparser ones; toward zero the
spacing shrinks down to the subnormal limit.}
\label{fig:titikmesin}
\end{figure}

There is one subtlety worth naming before we use the ulp as a unit of error. The largest
rounding error is half an ulp, but expressed as a relative error that figure is not one
number but an interval. For $x$ exactly equal to a power of two, half an ulp is
$2^{-53}$ relative; for $x$ that nearly reaches the next power of two, the value
approaches $2^{-54}$. So the relative error of one rounding can swing by a factor of
two. Goldberg named that swing \emph{wobble}, and that is why the bound $u$ is always
used as a worst case, not as a typical value.

\begin{teorema}[Relative rounding error]
For every $x\in\mathbb{R}$ in the normal range, correct rounding satisfies
$\operatorname{fl}(x)=x(1+\delta)$ with $|\delta|\le u=2^{-p}$.
\end{teorema}

\begin{proof}
Take $2^{e}\le|x|<2^{e+1}$. The nearest machine number is at most half an
ulp away, that is $2^{e-p}$, so that
$|\operatorname{fl}(x)-x|\le2^{e-p}=u\cdot2^{e}\le u|x|$. Dividing by $|x|$
gives $|\delta|\le u$, and the sign of $\delta$ follows the direction of rounding.
\end{proof}

The form $x(1+\delta)$ is the beginning of everything: it turns the question ``how
wrong'' into ``how many factors $(1+\delta)$ pile up''. The operations
$+,-,\times,\div,\sqrt{\phantom{x}}$ in IEEE-754 are required to return the correctly
rounded exact result, so every operation contributes one such factor
and no more. This guarantee is what makes error analysis possible.
\section{Error, condition number, and cancellation}
\tingkat{intermediate}

Now we have a unit. What we do not yet have is a measure of how much the error of a
computation exceeds the error of its input. Two problems can both receive input that is
wrong by $10^{-16}$ and come out with an answer that is wrong by $10^{-16}$ or by
$10^{-1}$, depending on the problem itself.

\kamus{condition number}{bilangan kondisi}{The multiplier showing how many times a
computation amplifies the error of its input.}\index{condition number}

\begin{definisi}[Absolute, relative, and backward error]
Let $\tilde{x}$ be an approximation to $x$. Its absolute error is $|\tilde{x}-x|$, its
relative error is $|\tilde{x}-x|/|x|$ for $x\neq0$. If $\tilde{x}$ is the exact answer
to a problem with data $x+\Delta x$, then $\|\Delta x\|/\|x\|$ is called the backward
error, and an algorithm that always gives a small backward error is called
\emph{backward stable}.\index{relative error}\index{backward stable}
\end{definisi}

The last distinction is the most useful in practice. A backward stable algorithm may be
badly wrong, provided its answer is right for a slightly different problem: the error
belongs to the problem. A large backward error is the mistake of the person writing the
code.

For a function of one variable, that amplification is computed directly with the
derivative: the condition number is the slope of the function at that point, multiplied
by the size of the point, divided by the size of its value.

\begin{teorema}[Condition number]
Let $f$ be differentiable on an interval containing $x$ and $f(x)\neq0$. If
$\tilde{x}=x(1+\delta)$ with $|\delta|\le u$, then
\[
  \frac{|f(\tilde{x})-f(x)|}{|f(x)|}\;\le\;\kappa_{f}(x)\,u+O(u^{2}),
  \qquad \kappa_{f}(x)=\frac{|x\,f'(x)|}{|f(x)|}.
\]
\end{teorema}

\begin{proof}
By the mean value theorem there is $\xi$ between $x$ and $\tilde{x}$ with
$f(\tilde{x})-f(x)=f'(\xi)(\tilde{x}-x)$. Since $|\tilde{x}-x|\le u|x|$ and $f'$ is
continuous at $x$, we have $|f'(\xi)|=|f'(x)|+O(u)$. Hence
$|f(\tilde{x})-f(x)|\le u|x|\,|f'(x)|+O(u^{2}|x|)$. Dividing by $|f(x)|$ gives the
claimed inequality.
\end{proof}

The number $\kappa_{f}$ does not depend on the machine, only on the problem. For
$f(x)=\sqrt{x}$ it is $1/2$: a square root never makes the relative error worse.
For $f(x)=\log x$ it is $1/|\log x|$, which is large when $x$ is near $1$ because
there $\log x$ is almost zero. For $f(x)=x-c$ with $c$ a large constant, it is
$|x|/|x-c|$, which grows without bound as $x$ approaches $c$. That last pattern is
what becomes the source of the whole trouble in this chapter.

Subtracting two numbers that are close together is the most dangerous operation in
machine arithmetic, and the danger is often blamed on the wrong thing. A common
misunderstanding runs: subtraction loses digits, so the operation of subtraction is
itself inaccurate. The truth is subtler and more reassuring: when the two numbers are
close together, the subtraction is in fact exact.

\begin{teorema}[Sterbenz]
If $a,b\in\mathbb{F}$ and $b/2\le a\le2b$ with $b>0$, then $a-b$ is represented exactly
in $\mathbb{F}$.
\end{teorema}

\begin{proof}
Let $2^{e}\le b<2^{e+1}$. Every machine number on that interval is an integer multiple of
$2^{e-p+1}$, so $a$ is as well, and hence $a-b$ is a multiple of $2^{e-p+1}$. From
$b/2\le a\le2b$ we get $0\le a-b\le b/2<2^{e}$. A multiple of $2^{e-p+1}$ that
lies in $[0,2^{e})$ has the form $M\cdot2^{e-p+1}$ with $0\le M<2^{p}$, and
all of them are represented with $p$ digits. Hence the result is exact.
\end{proof}

So when the result of a subtraction looks bad, what is bad is the input. The digits that
``disappear'' had already disappeared before the subtraction was carried out: $a$ and $b$
are each known only to relative error $u$, and an absolute error of size $u|a|$ does not
shrink along with the result when the two numbers cancel each other. Only the size of the
result shrinks, so that the relative error explodes.

\begin{proposisi}[Condition number of subtraction]
Let $x=a-b\neq0$ and let the machine work with $\tilde{a}=a(1+\delta_{a})$,
$\tilde{b}=b(1+\delta_{b})$, $|\delta_{a}|,|\delta_{b}|\le u$. Then
\[
  \frac{|\operatorname{fl}(\tilde{a}-\tilde{b})-x|}{|x|}
  \;\le\;\kappa\,u+u,
  \qquad \kappa=\frac{|a|+|b|}{|a-b|}.
\]
\end{proposisi}

\begin{proof}
Since $\operatorname{fl}(\tilde{a}-\tilde{b})$ is the correctly rounded version of
$\tilde{a}-\tilde{b}$, the additional error is at most $u|\tilde{a}-\tilde{b}|$.
Furthermore
$|\tilde{a}-\tilde{b}-(a-b)|\le|a||\delta_{a}|+|b||\delta_{b}|\le u(|a|+|b|)$,
and dividing by $|x|=|a-b|$ gives $u\kappa$. Adding the two errors with the
triangle inequality completes the proof.
\end{proof}

The number $\kappa$ is never less than one and grows without bound as
$|a-b|$ shrinks. For $a=12345.678$ and $b=12345.671$ we get $a-b=0.007$
and $\kappa\approx3.53\times10^{6}$, so that the relative error of the result can reach
$3.9\times10^{-10}$ even though each input is off by only $10^{-16}$. Since every
factor of ten in $\kappa$ erases one significant digit,
$\log_{10}\kappa\approx6.5$ means six and a half of the sixteen machine digits are gone,
and the result $0.007$ leaves only about ten digits that can be trusted.

One note that saves us: Sterbenz's rule applies to differences, not to
sums. For the sum $a+b$ the condition number is
$(|a|+|b|)/|a+b|$, which approaches one when $a$ and $b$ have the same sign. Adding
large numbers is not dangerous. Subtracting two large numbers that are nearly equal is
dangerous, and that is the only habit that needs watching.

\section{Catastrophic cancellation: two faces}
\tingkat{intermediate}

We already know the mechanism. Now let us give it a name, because two similar kinds of
cancellation are often taken for the same thing although their consequences differ
widely.

\begin{untukanak}
If we want to know the thickness of a sheet of paper, an ordinary ruler will not help:
one sheet is far thinner than a millimeter, so the needle seems not to move. The way is
not to look for a finer ruler but to change what is measured: stack five hundred sheets,
measure the thickness of the stack, then divide by five hundred. The result is far more
trustworthy, even though the ruler is the same ruler.

It turns out that a formula can be treated in the same way. Some formulas work by
subtracting two large numbers that are nearly equal, and that is the same as weighing a
single sheet of paper with a ruler. Other formulas are mathematically exactly the same,
but order their steps so that the small number is measured directly rather than
recovered from a difference. The results can differ widely, even though the paper has
not changed.

The question: if two formulas that are equally correct give different answers, which
one should we trust? In this chapter we shall learn why the wrong formula is not the one
whose arithmetic is mistaken but the one that chooses the step that throws information
away, and how to recognize it before the program runs.
\end{untukanak}

\kamus{catastrophic cancellation}{pembatalan katastrofik}{The loss of significant digits
when two nearly equal numbers are subtracted, so that only error is left.}
\index{catastrophic cancellation}

\begin{definisi}[Catastrophic and benign cancellation]
Let $x$ and $y$ be machine numbers computed from earlier data.
The subtraction $x-y$ is called \emph{benign cancellation} if $x$ and $y$ are exact, so
that the result is exact by Sterbenz's theorem. It is called \emph{catastrophic
cancellation} if $x$ and $y$ themselves already carry rounding error, so that the digits
the two share disappear as well and what is left is dominated by that error.
\end{definisi}

A benign example: $1.0000000000000002-1$ in binary64. Both are machine numbers, the
difference is $2^{-52}$ and is represented exactly. The result is small, but not one
digit disappears: we really do know both inputs to the last bit.

The catastrophic example is the quadratic formula from the opening of the chapter. There
$b^{2}$ and $4ac$ are each the result of a multiplication, so both are rounded. For
$a=1$, $b=10^{8}$, $c=1$ both happen to be exact, and it is precisely that specialness
which makes this example worth examining down to the bit: the collapse does not come
from the multiplications but from the rounding of the square root. Python computes

\begin{center}
\begin{tabular}{ll}
\hline
discriminant $b^{2}-4ac$ & \texttt{9999999999999996.0} \\
$\sqrt{b^{2}-4ac}$ & \texttt{99999999.99999999} \\
$b-\sqrt{b^{2}-4ac}$ & \texttt{1.4901161193847656e-08} \\
$\operatorname{ulp}(b)$ & \texttt{1.4901161193847656e-08} \\
\hline
\end{tabular}
\end{center}

The true value of $b-\sqrt{b^{2}-4ac}$ is $2\times10^{-8}$, that is $1.34$ ulp.
What is computed is exactly $1$ ulp, and half an ulp of error in the square root is
enough for that: the square root is correct to $0.46$ ulp at the scale $10^{8}$, but a
difference that ought to be $1.34$ ulp can only land on $1$ or $2$ ulp, so its
information is down to one binary choice. The final result is $-7.45\times10^{-9}$, that
is $-2^{-27}$, whereas the true root is $-10^{-8}$. Relative error $0.255$.

Note the division of labor. The square root works with relative error
$5\times10^{-17}$, far below $u$. The problem itself is well conditioned: shifting
$b$ by one ulp shifts the small root by one ulp, so
$\kappa\approx|b|/\sqrt{b^{2}-4ac}\approx1$. What is bad is the \emph{form} we have
chosen: the condition number of the subtraction at that step is
\[
  \kappa=\frac{|b|+\sqrt{b^{2}-4ac}}{\left|\sqrt{b^{2}-4ac}-|b|\right|}
  =\frac{2\times10^{8}}{1.49\times10^{-8}}\approx1.34\times10^{16},
\]
so that $\kappa u\approx1.49$: the relative error bound passes one, which means not a
single digit is guaranteed.

Figure~\ref{fig:katastrofik} separates two things that are often mixed up at that step:
the size of the numbers, and how many digits still mean something after the subtraction.

\begin{figure}[htbp]
\centering
\begin{tikzpicture}[every node/.style={font=\footnotesize}]
  \draw[ink,thick,->] (1.2,3.20) -- (10.6,3.20);
  \foreach \dec/\tk in {-9/{$10^{-9}$},-6/{$10^{-6}$},-3/{$10^{-3}$},0/{$10^{0}$},
                        3/{$10^{3}$},6/{$10^{6}$},9/{$10^{9}$}} {
    \draw[ink] ({1.4+(\dec+9)*0.5},3.26) -- ({1.4+(\dec+9)*0.5},3.14);
    \node[below] at ({1.4+(\dec+9)*0.5},3.12) {\tk};
  }
  \fill[accent] ({1.4+17*0.5},3.20) circle (2.4pt);
  \node[above=2pt] at ({1.4+17*0.5},3.24) {$b=10^{8}$};
  \fill[layerteal] ({1.4+1.173*0.5},3.20) circle (2.4pt);
  \node[above=2pt] at ({1.4+1.173*0.5},3.24) {$b-\sqrt{\Delta}$};
  \draw[<->,thick,dashed,layergrey] ({1.4+1.173*0.5},3.72) -- ({1.4+17*0.5},3.72);
  \node[above] at ({1.4+9.1*0.5},3.74) {sixteen orders of magnitude};
  \foreach \i in {0,...,16} {
    \fill[accent!28] ({1.4+\i*0.45},1.60) rectangle ++(0.45,0.40);
    \draw[ink!45] ({1.4+\i*0.45},1.60) rectangle ++(0.45,0.40);
  }
  \foreach \i in {0,...,15} {
    \fill[layergrey!16] ({1.4+\i*0.45},0.85) rectangle ++(0.45,0.40);
    \draw[ink!45] ({1.4+\i*0.45},0.85) rectangle ++(0.45,0.40);
  }
  \fill[accent!70] ({1.4+16*0.45},0.85) rectangle ++(0.45,0.40);
  \draw[ink!45] ({1.4+16*0.45},0.85) rectangle ++(0.45,0.40);
  \node[left] at (1.35,1.80) {$b$};
  \node[left] at (1.35,1.05) {$b-\sqrt{\Delta}$};
  \draw[{Latex}-{Latex},layergrey] (1.40,0.62) -- ({1.4+16*0.45},0.62);
  \node[below] at ({1.4+8*0.45},0.60) {sixteen identical digits, canceling out};
  \node[right] at ({1.4+17*0.45},1.05) {one digit};
\end{tikzpicture}
\caption{Top: $b$ and $\sqrt{b^{2}-4ac}$ on a log scale, sixteen orders of
magnitude away from their difference. Bottom: the boxes are significant digits, seventeen
go in and sixteen of them are the same so that they cancel out. What remains is one
digit, namely the difference $1.49\times10^{-8}$ itself; since half an ulp of error in
the square root is enough to shift the last digit, the small root comes out with a
relative error of $0.255$.}
\label{fig:katastrofik}
\end{figure}

The fix is not more precision but choosing which root to compute first.
The product of the two roots is $c/a$, so once the large root $r_{1}$ is known, the
small root is obtained by a division, not by a subtraction:
\[
  r_{1}=\frac{-b-\sqrt{b^{2}-4ac}}{2a},\qquad r_{2}=\frac{c}{a\,r_{1}}
  =\frac{2c}{-b-\sqrt{b^{2}-4ac}}.
\]
The second form never subtracts two numbers that are close together. For the example
above it gives $-10^{-8}$ exactly, that is the correctly rounded root. The same formula
with the signs arranged according to the sign of $b$ is used in almost every library
that solves quadratic equations, and the reason is stability, not tidiness.

The second face of catastrophic cancellation appears in a place that looks harmless:
the difference $1-\cos x$ for small $x$. For small $x$,
\[
  \frac{1-\cos x}{x^{2}}\longrightarrow\frac{1}{2},
\]
because $\cos x=1-x^{2}/2+x^{4}/24-\cdots$. The formula is right, and it is also the
most natural one to write. But $\cos x$ is stored with relative error $u$, and because
$|\cos x|\le1$, its absolute error is about $u$. After division by $x^{2}$, that error
becomes a relative error of about $2u/x^{2}$. For $x=10^{-8}$ that factor is $2.2$:
every digit is wiped out.

Here are the results of an actual run for ten values of $x$, compared with a reference
computed at 50 digits and with the Taylor series $0.5-x^{2}/24+x^{4}/720$; for
$x\le10^{-3}$ the two comparators agree to the last digit. The ``naive'' column is
$(1-\cos x)/x^{2}$, the ``identity'' column is
$\tfrac12\left(\sin(x/2)/(x/2)\right)^{2}$, which comes from
$1-\cos x=2\sin^{2}(x/2)$.

\begin{center}
\begin{tabular}{rccc}
\hline
$x$ & naive & relative error & identity \\
\hline
$10^{-1}$ & $0.499583472197418$ & $1.1\times10^{-14}$ & $0.499583472197423$ \\
$10^{-2}$ & $0.499995833347366$ & $2.9\times10^{-13}$ & $0.499995833347222$ \\
$10^{-3}$ & $0.499999958325503$ & $1.6\times10^{-11}$ & $0.499999958333335$ \\
$10^{-4}$ & $0.499999996961265$ & $5.2\times10^{-9}$ & $0.499999999583333$ \\
$10^{-5}$ & $0.500000041370185$ & $8.3\times10^{-8}$ & $0.499999999995833$ \\
$10^{-6}$ & $0.500044450291171$ & $8.9\times10^{-5}$ & $0.499999999999958$ \\
$10^{-7}$ & $0.499600361081320$ & $8.0\times10^{-4}$ & $0.500000000000000$ \\
$10^{-8}$ & $0.000000000000000$ & $1.0$ & $0.500000000000000$ \\
$10^{-9}$ & $0.000000000000000$ & $1.0$ & $0.500000000000000$ \\
$10^{-10}$ & $0.000000000000000$ & $1.0$ & $0.500000000000000$ \\
\hline
\end{tabular}
\end{center}

The pattern is typical and worth memorizing. As long as $x^{2}/2$ is far above $u$, the
naive formula is nearly right. When $x^{2}/2$ shrinks to the scale of $u$, the
significant digits run out and the error grows like $2u/x^{2}$: each tenfold decrease of
$x$ adds two digits of loss. At $x=10^{-8}$, $1-\cos x$ is smaller than half an ulp
around $1$, so that $\cos x$ is rounded exactly to $1$ and the result is zero. That
breakdown point is no mystery: $x^{2}/2\approx u$, that is $x\approx\sqrt{2u}\approx1.5\times10^{-8}$.

The last column shows that the problem is well conditioned. The double-angle identity
gives an answer correct to the last digit for every $x$ in the table, because it never
subtracts two numbers that are close together. Error is not a property of the problem.
It is a property of the formula chosen.

\begin{resep}
Four steps for choosing a formula form that is safe from cancellation:
\begin{enumerate}
  \item Mark every subtraction $a-b$ inside the formula, then compute
        $\kappa=(|a|+|b|)/|a-b|$ on the worst data that can reasonably arrive.
  \item Translate $\kappa$ into lost digits: $\kappa u\approx10^{-d}$
        means the last $d$ digits cannot be trusted.
  \item Look for an equivalent form that contains no such difference: multiply by the
        conjugate, use the double-angle identity, factor, or call \texttt{log1p},
        \texttt{expm1}, and \texttt{hypot}.
  \item Repeat the computation at higher precision. The difference between the two
        results should be comparable to $\kappa u$; if it is far larger, the new form
        still hides a subtraction.
\end{enumerate}
A warning about order: step 1 must not be swapped with step 3. Rewriting a formula
before knowing which subtraction is dangerous only moves the problem, as in the one-pass
form of the variance that trades one dangerous subtraction for another.
\end{resep}
\section{Summation: error that grows with $n$}
\tingkat{advanced}

Summation is the operation most often taken for granted, and yet it is where error is
easiest to compute. Before we come to the bound, here is a demonstration that dismantles
the assumption that order does not matter:

\begin{center}
\texttt{(0.1 + 0.2) + 0.3 = 0.6000000000000001}, \quad
\texttt{0.1 + (0.2 + 0.3) = 0.6}.
\end{center}

The difference is $1.11\times10^{-16}$, exactly one ulp at the scale $0.6$. No digit is
written down wrongly: only three numbers, each with no finite binary
representative, summed in a different order. In mathematics, addition is
associative. Inside the machine, it is associative only when no differing rounding
happens to occur.

\begin{teorema}[Error of sequential summation]
Let $s_{1}=x_{1}$ and $s_{k}=\operatorname{fl}(s_{k-1}+x_{k})$ for
$k=2,\dots,n$ with $x_{i}\in\mathbb{F}$. Then
\[
  \left|s_{n}-\sum_{i=1}^{n}x_{i}\right|
  \;\le\;(n-1)\,u\sum_{i=1}^{n}|x_{i}|+O(u^{2}).
\]
\end{teorema}

\begin{proof}
Every step uses one rounding, so $s_{k}=(s_{k-1}+x_{k})(1+\delta_{k})$
with $|\delta_{k}|\le u$. Expanding the result recursively gives
\[
  s_{n}=\sum_{i=1}^{n}x_{i}\prod_{k=i}^{n}(1+\delta_{k}),
\]
where the first term carries no $\delta_{1}$ because $s_{1}=x_{1}$ has not been rounded.
Since $\prod(1+\delta_{k})=1+\theta$ with
$|\theta|\le(1+u)^{n-1}-1\le(n-1)u+O(u^{2})$, the triangle inequality gives the bound
claimed.
\end{proof}

That bound is honest but loose: it assumes every rounding reaches its
worst error and that all of them point the same way, which in practice never happens.
The following table shows the gulf between bound and reality for the harmonic series
$H_{n}=\sum_{i=1}^{n}1/i$ in binary32, with the exact value computed by
\texttt{math.fsum}, which sums exactly.

\begin{center}
\begin{tabular}{rccc}
\hline
$n$ & sequential & exact & bound $(n-1)u\sum|x_{i}|$ \\
\hline
$10^{2}$  & $5.1873779297$ & $5.1873775176$ & $3.1\times10^{-5}$ \\
$10^{3}$  & $7.4854784012$ & $7.4854708606$ & $4.5\times10^{-4}$ \\
$10^{4}$  & $9.7876129150$ & $9.7876060360$ & $5.8\times10^{-3}$ \\
$10^{5}$  & $12.0908508301$ & $12.0901461299$ & $7.2\times10^{-2}$ \\
$10^{6}$  & $14.3573579788$ & $14.3927267229$ & $8.6\times10^{-1}$ \\
\hline
\end{tabular}
\end{center}

The actual error at $n=10^{6}$ is $3.54\times10^{-2}$, about twenty-five
times smaller than the bound, but it has already erased two and a half significant
digits. Nor does the error rise smoothly: it dips at $n=10^{4}$ because roundings in the
same and in opposite directions happened to cancel. What can be relied on is not the
size of the error at one $n$, but its tendency to grow, at a worst-case rate
proportional to $n$.

The cause is not large rounding but an accumulator that moves away from zero.
Once $s$ is around $14$, the term $1/i$ for large $i$ lies below half an
ulp of $s$ and every addition loses all its information. The extreme case is easy to
construct: add $1$ to a million terms of value $10^{-16}$; exactly the result is
$1.0000000001$, whereas sequential summation stays at $1$ for however many terms there
are. What is lost is not one term but the whole contribution.

Figure~\ref{fig:galatjumlah} summarizes the same pattern for the five values of $n$ in the
table above, this time with Kahan summation as the comparison.

\begin{figure}[htbp]
\centering
\begin{tikzpicture}
\begin{axis}[
  width=0.88\linewidth, height=6.0cm,
  xmode=log, ymode=log,
  xlabel={$n$}, ylabel={$|s_{n}-H_{n}|$},
  xmin=70, xmax=1.6e6, ymin=1e-8, ymax=1,
  grid=both, grid style={layergrey!15},
  tick label style={font=\footnotesize},
  label style={font=\footnotesize},
  legend style={font=\footnotesize,draw=layergrey!40,
                at={(0.03,0.97)},anchor=north west},
  legend cell align=left,
]
\addplot[ink,thick,mark=*,mark size=1.7pt] coordinates {
  (100,4.12e-7) (1000,7.54e-6) (10000,6.88e-6) (100000,7.05e-4) (1000000,3.54e-2)};
\addlegendentry{sequential}
\addplot[accent,thick,mark=square*,mark size=1.6pt] coordinates {
  (100,6.48e-8) (1000,8.88e-8) (10000,2.03e-7) (100000,6.51e-8) (1000000,1.75e-7)};
\addlegendentry{Kahan}
\end{axis}
\end{tikzpicture}
\caption{Absolute error of the harmonic sum in binary32 at five sizes $n$, log-log
scale. Sequential summation rises from $4.1\times10^{-7}$ at $n=10^{2}$ to
$3.54\times10^{-2}$ at $n=10^{6}$; Kahan holds at order $10^{-7}$ across the whole
range, so that at $n=10^{6}$ the two are $2\times10^{5}$ times apart. The data, the
precision, and the number of operations are the same; only the order of summation
differs.}
\label{fig:galatjumlah}
\end{figure}

\begin{tangan}
Materials: eleven blank cards, a marker, and one rule of the game. Every time two cards
are added, write the result on a new card with at most four significant figures,
then throw the old cards away.

Write $1.000$ on one card and $0.0001$ on ten other cards. First route: put down the big
card, then add the ten small cards one at a time. Each time
$1.000+0.0001=1.0001$ is cut down to $1.000$, so that the small card disappears
without a trace and the pile stays $1.000$.

Shuffle the cards again, second route: first add up the ten small cards into
$0.0010$, then add the big card. The final pile is $1.001$.

What must be visible: the same eleven cards, the same cutting rule, two
different results in the last digit. Only the order changed, and the machine repeats that
order a million times.
\end{tangan}

Kahan summation attacks that problem directly.\index{Kahan summation}
The idea is simple: every time
a term is lost to rounding, keep the lost remainder in a compensation variable,
then put it back in at the next step.

\begin{quote}
\texttt{s = 0.0; c = 0.0}\\
\texttt{for x in data:}\\
\hspace*{1.5em}\texttt{y = x - c}\\
\hspace*{1.5em}\texttt{t = s + y}\\
\hspace*{1.5em}\texttt{c = (t - s) - y}\\
\hspace*{1.5em}\texttt{s = t}
\end{quote}

The line \texttt{c = (t - s) - y} is the heart of it, and it works precisely because of
Sterbenz: $t-s$ is the difference of two numbers that are close together and hence exact,
and $(t-s)-y$ is therefore exactly the error that was just discarded. The result no
longer piles up error proportional to $n$.

\begin{teorema}[Kahan summation bound]
With Kahan's algorithm above, the computed sum satisfies the inequality
\[
  s=\sum_{i=1}^{n}x_{i}(1+\delta_{i})+O(nu^{2})\sum_{i=1}^{n}|x_{i}|,
  \qquad |\delta_{i}|\le2u .
\]
\end{teorema}

We do not repeat the proof of this theorem here. It is Theorem 8 in Goldberg's
survey, and its proof is in the appendix of that work, taken from Knuth. What is worth
reading from the statement: the relative error of each term is bounded by $2u$, twice the
error of a single rounding, but the factor $n$ is gone from the leading term. For small
$n$ this bound is in fact looser than the naive one, because $2u>u$. Its advantage
appears only when $n$ is large, and that is where its whole use lies.

The real numbers on the harmonic series with $n=10^{6}$: sequential summation in
binary32 gives an error of $3.54\times10^{-2}$, whereas Kahan gives
$1.75\times10^{-7}$, that is $0.18$ ulp from the correct answer. In binary64
the difference is just as sharp: sequential is off by $7.35\times10^{-13}$, or $414$ ulp,
whereas Kahan produces a figure identical to \texttt{math.fsum}.

Two other improvements deserve mention. Pairwise summation, which
\texttt{numpy.sum} uses,\index{pairwise summation} splits the data recursively
so that the error grows like $\log n$; it is nearly free and suits data in
memory. \texttt{math.fsum} uses an exact accumulator and gives a result that is correctly
rounded for the whole list, at a higher cost in time. Summing from the smallest term to
the largest is the cheapest improvement, but it only reduces the error.

\section{Practical guidance}
\tingkat{basic}

Everything we have built can be compressed into a few habits that are cheap to
keep and expensive to neglect.

First, choose precision from the condition number, not from habit.
The relative error of the result is roughly $\kappa u$. With binary64, a computation
with condition $10^{6}$ still leaves ten digits; with binary32, the same computation
leaves only one digit. For $10^{8}$ sequential operations the crude bound is
$nu$, and in binary32 that already passes one, which means the answer is meaningless.

Second, do not test fractional numbers for equality with \texttt{==}. Two paths that
are mathematically identical may produce different bits, like
$0.1+0.2$ above. The right test uses a scaled tolerance, for instance
$|x-y|\le10^{-12}\max(1,|x|,|y|)$, and that tolerance must come from error
analysis, not from experimenting until the test passes.

Third, distrust every formula that contains a difference of two nearly equal quantities.
The standard form of the variance, $(1-\cos x)/x^{2}$, and $\log(1+x)$ for small $x$ can
all be written differently. The standard library already does so: \texttt{log1p},
\texttt{expm1}, and \texttt{hypot} exist precisely because their naive forms are not
stable.

Fourth, for a long sum, or one with terms of very different sizes, use
\texttt{math.fsum} or pairwise summation, and keep the accumulator in higher precision
where available. Sort from the smallest term if there is no other choice.

Fifth, do not conclude truth from agreement between numbers. Two programs that agree
to twelve digits are fairly convincing; ones that agree to three digits only tell
us that both use the same formula. A better test is to repeat the computation at a
different precision: if the result shifts far beyond the prediction
$\kappa u$, something is wrong with the formula, not with the machine.


\pagebreak[4]
\begin{tantangan}
Compute $10^{8}\left(\sqrt{10^{16}+1}-10^{8}\right)$ in binary64. The true value is
$0.5$, but a program that writes the formula down as it stands reports $0.0$.

Explain why the result is zero, then compute the form that gives $0.5$ without raising
the precision at all.

Hint 1: ask what $\operatorname{ulp}(10^{16})$ is, then whether $10^{16}+1$
still has room on the grid of machine numbers.

Hint 2: multiply numerator and denominator by the conjugate
$\sqrt{10^{16}+1}+10^{8}$; the dangerous subtraction becomes a division.

If you are stuck, go back to the section Catastrophic cancellation: two faces.
\end{tantangan}

\section{Worked example}
\tingkat{intermediate}

Every number in this section can be checked with a calculator, \texttt{struct}, and
ordinary fractions. We begin by taking apart one bit pattern, namely $0.1$ in binary32:

\[
  \underbrace{0}_{1}\;\underbrace{01111011}_{8}\;
  \underbrace{10011001100110011001101}_{23}.
\]

The stored exponent is $123$, so after subtracting the bias $127$ we get $e=-4$. The
23-bit integer in the mantissa is $5033165$, so that the significand is
$1+5033165/2^{23}$. Hence
\[
  x=\left(1+\frac{5033165}{2^{23}}\right)2^{-4}=\frac{13421773}{134217728}.
\]
The numerator comes from $2^{23}+5033165=13421773$. Meanwhile the denominator is
$134217728$, and a tenth of it is $13421772.8$, so that the error is exactly
$0.2/134217728$, that is $1.4901161\times10^{-9}$,
and the relative error is $1.49\times10^{-8}$, that is $u_{32}/4$. That figure makes
sense: the number $1/10$ equals $1.1001100110011\dots_{2}\times2^{-4}$, which needs a
periodic mantissa and does not terminate, so that correct rounding cuts it at the 24th
bit. The pattern that is stored
ends in \texttt{\dots 1101}, not \texttt{\dots 1100}, because the remainder behind it is
greater than half an ulp and the rounding goes up.

\kertas{Take apart the pattern \texttt{0 01111100 01000000000000000000000} by hand:
determine the sign, the exponent after subtracting the bias, the value of the mantissa,
then write the number as a fraction and as a decimal. Compute the
$\operatorname{ulp}$ of that number, then decide the result of adding $2^{-27}$ and
$3\cdot2^{-27}$ under the \emph{tie to even} rule. Check your answer with
\texttt{struct.unpack} and \texttt{numpy.float32}. This number was chosen so that its
mantissa stops after two digits, so that all of its arithmetic can be done without a
calculator.}

Now the case $0.1+0.2$ in binary64, which can be computed entirely with
fractions. All three numbers involved have power-of-two denominators:
\[
  \operatorname{fl}(0.1)=\frac{3602879701896397}{2^{55}},\qquad
  \operatorname{fl}(0.2)=\frac{3602879701896397}{2^{54}},
\]
\[
  \operatorname{fl}(0.3)=\frac{5404319552844595}{2^{54}}.
\]
The sum of the first two before rounding is $10808639105689191/2^{55}$,
whereas $\operatorname{fl}(0.3)=10808639105689190/2^{55}$. So the exact sum
lies exactly half an ulp above $\operatorname{fl}(0.3)$: a tie point, not something near
one. The last digit of its mantissa is $1$, that is odd, so that \emph{tie to even}
rounds up: $0.300000000000000044408920985006\dots$

The difference from $\operatorname{fl}(0.3)$ is therefore exactly one ulp,
$2^{-54}=5.551115123125783\times10^{-17}$, not a random error. The subtraction that
follows is exact by Sterbenz, so that the figure the program prints is that figure
precisely.

The same numbers explain why the conclusion depends on the precision: in
binary32, $0.1+0.2$ and $0.3$ produce identical bit patterns because the two
rounding errors happen to cancel each other. So \texttt{0.1+0.2 != 0.3} is not a law of
nature but the consequence of a coincidence at one precision.

Now the square root from before, with ulp arithmetic. The discriminant $b^{2}-4ac$ is
exactly $10^{16}-4$, because $10^{16}$ is still below $2^{54}$ and that number is a
multiple of $4$. Its square root is correct to $0.46$ ulp, so that it lands one ulp
below $b$ because $\operatorname{ulp}(b)=2^{-26}$. The true difference is $2\times10^{-8}=1.34$ ulp,
the computed difference $1$ ulp; that error of $0.34$ ulp at the scale $10^{8}$, divided
by two, becomes the whole error of the small root:
\[
  r_{2}^{\text{naive}}=-\frac{2^{-26}}{2}=-2^{-27}
  =-7.450580596923828\times10^{-9},
  \qquad r_{2}=-10^{-8}.
\]
The stable form $2c/(-b-\sqrt{b^{2}-4ac})$ suffers no dangerous subtraction, and
its result is $-10^{-8}$ exactly, the correctly rounded root.
\section{Computational lab}
\tingkat{advanced}

\begin{lab}
\begin{verbatim}
import math, struct
import numpy as np

def bit32(x):
    return ''.join(f'{b:08b}' for b in struct.pack('>f', float(x)))

print("0.1     :", bit32(0.1))
print("6.0     :", bit32(6.0))
print("-12.375 :", bit32(-12.375))
print("hex     :", (0.1).hex(), (6.0).hex(), (1e-8).hex())
print("eps     :", np.finfo(np.float32).eps, np.finfo(np.float64).eps)
print("ulp     :", math.ulp(1.0), math.ulp(1e10))
print("tie     :", 1 + 2.0**-53, 1 + 3*2.0**-53)

s = (0.1 + 0.2) - 0.3
print("0.1+0.2-0.3 :", repr(s), "=", s / math.ulp(0.3), "ulp")
print("associative :", repr((0.1 + 0.2) + 0.3), repr(0.1 + (0.2 + 0.3)))

n = 10**6
exact = math.fsum(1.0 / i for i in range(1, n + 1))

def jumlah(prec, kahan):
    s, c = prec(0), prec(0)
    for i in range(1, n + 1):
        y = prec(1) / prec(i)
        if kahan:
            y = y - c
            t = s + y
            c = (t - s) - y
            s = t
        else:
            s = s + y
    return s

print("H_n exact   :", f"{exact:.12f}")
for prec, nama in ((np.float32, "float32"), (np.float64, "float64")):
    for kahan in (False, True):
        v = float(jumlah(prec, kahan))
        print(f"{nama} {'Kahan' if kahan else 'naive':7} {v:.12f}"
              f"  error {abs(v - exact):.2e}")

print(f"{'x':>8} {'naive':>19} {'error':>9} {'identity':>19} {'error':>9}")
for k in range(1, 11):
    x = 10.0**-k
    lugu = (1 - math.cos(x)) / x**2
    stabil = 2 * math.sin(x / 2)**2 / x**2
    acuan = 0.5 - x**2/24 + x**4/720 - x**6/40320
    print(f"{x:8.0e} {lugu:19.15f} {abs(lugu-acuan)/acuan:9.1e}"
          f" {stabil:19.15f} {abs(stabil-acuan)/acuan:9.1e}")

a, b, c = 1.0, 1e8, 1.0
d = b*b - 4*a*c
sq = math.sqrt(d)
print("discriminant :", d, " root:", sq)
print("b - root    :", b - sq, " ulp(b):", math.ulp(b))
print("naive root  :", (-b + sq) / (2*a))
print("stable root :", 2*c / (-b - sq))
\end{verbatim}
The actual output:
\begin{verbatim}
0.1     : 00111101110011001100110011001101
6.0     : 01000000110000000000000000000000
-12.375 : 11000001010001100000000000000000
hex     : 0x1.999999999999ap-4 0x1.8000000000000p+2 0x1.5798ee2308c3ap-27
eps     : 1.1920929e-07 2.220446049250313e-16
ulp     : 2.220446049250313e-16 1.9073486328125e-06
tie     : 1.0 1.0000000000000004
0.1+0.2-0.3 : 5.551115123125783e-17 = 1.0 ulp
associative : 0.6000000000000001 0.6
H_n exact   : 14.392726722866
float32 naive   14.357357978821  error 3.54e-02
float32 Kahan   14.392726898193  error 1.75e-07
float64 naive   14.392726722865  error 7.35e-13
float64 Kahan   14.392726722866  error 0.00e+00
       x               naive     error            identity     error
   1e-01   0.499583472197418   5.6e-15   0.499583472197423   5.4e-15
   1e-02   0.499995833347366   2.9e-13   0.499995833347222   1.1e-16
   1e-03   0.499999958325503   1.6e-11   0.499999958333335   0.0e+00
   1e-04   0.499999996961265   5.2e-09   0.499999999583333   1.1e-16
   1e-05   0.500000041370185   8.3e-08   0.499999999995833   1.1e-16
   1e-06   0.500044450291171   8.9e-05   0.499999999999958   0.0e+00
   1e-07   0.499600361081320   8.0e-04   0.500000000000000   1.1e-16
   1e-08   0.000000000000000   1.0e+00   0.500000000000000   0.0e+00
   1e-09   0.000000000000000   1.0e+00   0.500000000000000   0.0e+00
   1e-10   0.000000000000000   1.0e+00   0.500000000000000   0.0e+00
discriminant : 9999999999999996.0  root: 99999999.99999999
b - root    : 1.4901161193847656e-08  ulp(b): 1.4901161193847656e-08
naive root  : -7.450580596923828e-09
stable root : -1e-08
\end{verbatim}
\end{lab}

The first three lines confirm that the table at the beginning of the chapter is not
invented: the bit patterns of $0.1$, $6$, and $-12.375$ come out exactly as in the hand
decomposition, sign bit included. The \texttt{hex} line shows thirteen hexadecimal
digits. The \texttt{tie} line is the plainest demonstration of the even rule: adding
half an ulp to $1$ changes nothing, adding one and a half ulp gives
$1+2^{-51}$. The \texttt{0.1+0.2-0.3} line prints
$5.551115123125783\times10^{-17}$ and declares itself to be $1.0$ ulp; the
\texttt{associative} line gives two different results, differing by
$1.11\times10^{-16}$.

The next four lines are the core of this chapter in its most compact form: at
$n=10^{6}$, binary32 sequential loses $3.54\times10^{-2}$ whereas Kahan loses only
$1.75\times10^{-7}$; binary64 sequential $7.35\times10^{-13}$ whereas Kahan zero, the
same bits as \texttt{math.fsum}. That error of $7.35\times10^{-13}$ is equivalent to
$414$ ulp, so double precision only moves the problem to a smaller scale.

The table for $(1-\cos x)/x^{2}$ shows the pattern already computed on paper. At
$x=10^{-1}$ both methods differ from the reference by about $5\times10^{-15}$ because the
reference itself is truncated at the $x^{6}$ term; from $x=10^{-2}$ on that reference is
far more accurate than the machine, and the error of the naive formula rises from
$2.9\times10^{-13}$ to $1.0$ as $x$ falls to $10^{-8}$, whereas the identity column holds
at $10^{-16}$. The zero at $x=10^{-8}$ is not lazy rounding: $\cos(10^{-8})$ really
equals $1$.

The last block checks the square root: discriminant $10^{16}-4$, root $10^{8}-2^{-26}$,
a difference of one ulp of $b$, namely $2^{-26}$, naive root $-2^{-27}$, and stable root
$10^{-8}$, which is correctly rounded. The last two lines
are $25.5$ percent apart, even though the formulas are algebraically equivalent and the
data are exactly the same.

\begin{untukdosen}
A teaching order that has worked in practice: begin with one failure that can be
recomputed, namely the quadratic formula at the opening of the chapter, and only then go
into representation. If representation comes first, students memorize $1/8/23$ and
$1/11/52$ without knowing what for. After that failure, the order that works is: the
number of machine numbers (so that it is clear rounding cannot be avoided), the rounding
rule, the ulp lemma, the $(1+\delta)$ bound, then the condition number. Summation is
better placed after the condition number, because the bound $(n-1)u\sum|x_{i}|$ only
reads as ``error that grows'' once $\kappa$ is familiar.

Two typical misconceptions. First, students believe that rounding error is random and
therefore cancels out in long sums. The harmonic table contradicts that sharply: the
error moves in one direction, and at $n=10^{6}$ it reaches thirty-five thousand ulp.
Second, they believe that subtraction always damages digits, and then avoid all
subtraction. Sterbenz's theorem shows the opposite: the subtraction of two nearby machine
numbers is in fact exact, and what does the damage is the error already present in the
input. The rule that needs repeating: what is dangerous is not the minus sign, but a
small difference of two numbers that have each already been rounded.

A cheap and effective whiteboard activity: ask the whole class to compute
$1+2^{-53}$ and $1+3\cdot2^{-53}$ with the even rule, then compare with Python.
After that, ask two volunteers to work out
$(0.1+0.2)+0.3$ and $0.1+(0.2+0.3)$ by writing the 52-bit mantissa
as a fraction. The hubbub usually starts when they discover that both are right.

A sign of readiness to move on: students can guess, without running a program, that
$\operatorname{fl}(1+10^{-16})=1$, can explain why the form
$2c/(-b-\sqrt{\Delta})$ is better than $(-b+\sqrt{\Delta})/(2a)$, and can
estimate how many digits are lost from a given $\kappa$. If
they say the words ``backward error'' without being prompted, they are ready for Chapter 25,
where the gradient of a network is computed with a similar rule.
\end{untukdosen}

\section{Summary}
\tingkat{basic}

\begin{ringkasan}
\begin{itemize}
  \item $\mathbb{F}$ is finite: $2^{p-1}$ normal numbers per power of two, so
        almost every real number has no representative and every act of storage means
        rounding.
  \item IEEE rounding is \emph{round to nearest, ties to even}. Consequently
        $\operatorname{fl}(x)=x(1+\delta)$ with $|\delta|\le u=2^{-p}$, that is
        $1.11\times10^{-16}$ for binary64 and $5.96\times10^{-8}$ for
        binary32.
  \item $\operatorname{eps}=2^{-p+1}$ is the distance from $1$ to the number after it;
        $\operatorname{ulp}(x)=2^{e-p+1}$ for $2^{e}\le|x|<2^{e+1}$. The error of one
        rounding is half an ulp, and expressed as a relative error it swings by
        a factor of two.
  \item The condition number $\kappa_{f}=|xf'(x)|/|f(x)|$ measures the amplification of
        input error. For subtraction, $\kappa=(|a|+|b|)/|a-b|$, which grows without
        bound as $|a-b|$ shrinks.
  \item Subtracting two nearby machine numbers is exact (Sterbenz). Catastrophic
        cancellation happens because the inputs already carry error, not because of
        the subtraction.
  \item The naive quadratic formula computes the small root through a subtraction of two
        numbers of size $|b|$; the form $2c/(-b-\sqrt{\Delta})$ avoids it and gives the
        correctly rounded root.
  \item Sequential summation has a worst-case error of $(n-1)u\sum|x_{i}|$ and is not
        associative. Kahan moves the error into a compensation variable so that the
        factor $n$ disappears from the leading term.
\end{itemize}
\end{ringkasan}

\section{Exercises}
\tingkat{intermediate}

\begin{latihan}
\begin{enumerate}
  \item Determine the $\operatorname{ulp}$ of $1$; $1024$; $0.1$; and $3\times10^{8}$
        in binary64, then compare with $|x|\,u$. How many binary64
        numbers lie in $[1,2)$?

  \item Show that $1+2^{-53}=1$ and $1+2^{-52}>1$ in binary64, then determine
        the smallest positive number $x$ with $1+x>1$ in binary32. Express
        your answer as $2^{-k}$ and compare it with the table in the opening section.

  \item Prove that for every normal $x$ we have
        $|\operatorname{fl}(x)-x|\le\tfrac12\operatorname{ulp}(x)$, then explain
        why that inequality cannot be replaced by equality for all $x$.

  \item Assess this claim: \emph{because every IEEE-754 operation is correctly rounded,
        the relative error of the final result of a program that uses only $+,-,\times,\div$ is always
        at most a few $u$.} Decide whether it is true, false, or not supported by
        evidence, and name one example from this chapter that supports your decision.

  \item Compute by hand $(0.1+0.2)+0.3$ and $0.1+(0.2+0.3)$ as
        fractions with power-of-two denominators, then show that the difference is one ulp.
        Check with \texttt{math.fsum} and with \texttt{float.hex()}.
        (Hint: write each number as $m/2^{55}$ or $m/2^{54}$ and
        compute the numerators as integers.)

  \item Repeat the computation of $(1-\cos x)/x^{2}$ for $x=10^{-7}$ with the double-angle
        identity, then compare with $0.5-x^{2}/24+x^{4}/720$. What relative error
        remains, and where does it come from?

  \item For $a=1$, $b=10^{4}$, $c=1$, compute both roots with the naive formula and with
        the stable form in binary64. Determine the condition number of the subtraction, then
        predict how many digits are lost before running the program. Check
        your prediction with Python.
        (Hint: $\sqrt{10^{8}-4}\approx10^{4}-2\times10^{-4}$; compare with
        $\operatorname{ulp}(10^{4})=2^{-39}$.)

  \item Prove that pairwise summation over $n=2^{k}$ terms has a worst-case error
        proportional to $ku\sum|x_{i}|$, not $nu\sum|x_{i}|$, by expanding
        the summation tree and counting the depth of the tree.

  \item Let $x_{i}=10^{-16}$ for $i=1,\dots,10^{6}$ and $x_{0}=1$. Compute
        sequential summation and Kahan summation both by hand (a calculator may be
        used), then explain why the Kahan answer differs even though
        each of its steps is also rounded.
        (Hint: in sequential summation every term lies below half an ulp
        of the accumulator; in Kahan, the compensation $c$ keeps that remainder and adds it
        back at the next step.)

  \item A program computes the variance with
        $\frac1n\sum x_{i}^{2}-\left(\frac1n\sum x_{i}\right)^{2}$ on data that
        are all around $999$. Determine the condition number of that form, predict
        how many digits are lost in binary32, and compare with the
        two-pass form $\frac1n\sum(x_{i}-\bar{x})^{2}$. Explain why the second form is
        better even though it requires two passes over the data.
\end{enumerate}
\end{latihan}

\begin{untukpeneliti}
What is really proved in this chapter are standard statements that do not
depend on any particular machine: the bound $\operatorname{fl}(x)=x(1+\delta)$ with
$|\delta|\le u$, the ulp lemma, Sterbenz's theorem, the error bound for sequential
summation, and the condition-number inequality via the mean value theorem. The primary
reference is Goldberg,
\emph{What Every Computer Scientist Should Know About Floating-Point Arithmetic}
(ACM Computing Surveys 23(1), 1991), sections ``Relative Error and Ulps'',
``Cancellation'', ``Exactly Rounded Operations'', and ``Errors in Summation''. The Kahan
bound is Theorem 8 of that work; its proof is in the appendix, taken from Knuth, and is
not repeated here. For a systematic treatment of stability, the further reference is
Higham, \emph{Accuracy and Stability of Numerical Algorithms}, the chapters on error
analysis and summation.

Three things are more often popularized than proved. First, the claim that
``rounding error is random and cancels out''. That is true only as a statistical
model for uncorrelated rounding, and it fails precisely in the most important case:
accumulation in long sums is systematic, as the harmonic table shows. Second, the claim
that the error can always be reduced with higher precision. For catastrophic
cancellation, raising the precision only moves the breakdown point, it does not remove
it: in binary64 that point lies at
$x\approx1.5\times10^{-8}$, and in binary128 it lies around $10^{-17}$. Third, the claim
that the ulp is an unambiguous unit of error. Goldberg shows that the
ulp may be defined as the distance between numbers or as half that distance, and the
difference between the two definitions appears as a factor of two in every bound. The
figure ``$0.46$ ulp'' in the discussion of the square root in this chapter uses the first
definition.

The limits of validity of the results here must be stated. The whole analysis assumes a model of
correct rounding and no \emph{fused multiply-add} that changes the order of operations,
no rearrangement by the compiler, and no rounding mode other than
\emph{to nearest}. On real hardware, vectorization and FMA instructions can change
results bit for bit; that is not a violation of the standard, but it makes comparisons
between machines non-deterministic. Error analysis also does not touch model error,
data error, or truncation error: this chapter discusses only rounding error.

Honest open questions. First, classical error analysis gives worst-case bounds,
whereas what practitioners need is a tight estimate for the typical case;
the probabilistic treatment of rounding is still an actively researched tool and not a
finished theory. Second, for iterative algorithms and the solution of linear systems,
rounding error interacts with the condition number of the matrix in such a way that
naive bounds are far too pessimistic; backward analysis, which is the central idea of
Higham and Wilkinson, gives a sharper framework but is not discussed here.
Third, practical decisions such as ``is binary32 enough'' depend on the distribution of
the data, and no answer holds in general.
\end{untukpeneliti}

\peta{This chapter follows Goldberg, \emph{What Every Computer Scientist Should Know About
Floating-Point Arithmetic} (ACM Computing Surveys, 1991), sections ``Introduction'',
``Relative Error and Ulps'', ``Cancellation'', ``Exactly Rounded Operations'',
``Formats and Operations'', ``Special Quantities'', and ``Errors in Summation'', with
Theorem 8 for the Kahan summation bound. For the real-analysis side of the limit ideas
and the mean value theorem used in the proof of the condition number, see Abbott,
\emph{Understanding Analysis}, sections 2.2 (Limit of a Sequence), 5.3 (The Mean Value
Theorem), and 6.6 (Taylor Series). The example $0.1+0.2$, the decomposition of bit
patterns, the analysis of $(1-\cos x)/x^{2}$, and all the numbers in the tables and the
lab are computations of our own run with Python 3.14 and NumPy 2.5 and re-verified with
ordinary fractions and 50-digit arithmetic. For an advanced treatment of stability
and backward error, the reference commonly used is Higham, \emph{Accuracy and
Stability of Numerical Algorithms}.}
